% Report for Fig 2b/2d: Power Efficiency and Instantaneous Power Analysis
% To include in LaTeX: % Report for Fig 2b/2d: Power Efficiency and Instantaneous Power Analysis
% To include in LaTeX: % Report for Fig 2b/2d: Power Efficiency and Instantaneous Power Analysis
% To include in LaTeX: % Report for Fig 2b/2d: Power Efficiency and Instantaneous Power Analysis
% To include in LaTeX: \input{fig2b_power_efficiency.tex}

Fig.~\ref{fig:power_efficiency} shows that our architecture achieves 5.6--6.6$\times$ higher energy efficiency (Token/J) than H100 at BS=1, consistently across Qwen3-235B and Llama4-Maverick.
This advantage holds across batch sizes (BS=1--32) and sequence lengths (4K--1M), driven by the 40\,TB/s aggregate HBM bandwidth of 16~NPUs versus 3.35\,TB/s on H100.
Compared to Rubin, our design maintains a stable 2.6--3.1$\times$ Token/J advantage; against Rubin~TP2, the gap ranges from 2.8$\times$ (1M) to 4.8$\times$ (4K), as TP2's doubled static power and NVLink overhead erode its latency benefit at short sequences.

As shown in Fig.~\ref{fig:power_breakdown}, although our architecture draws higher instantaneous power than H100 (e.g., 1\,316\,W vs.\ 603\,W at BS=4, seq=4K on Llama4), it completes each token 12--16$\times$ faster, resulting in 82--85\% lower energy per token.
In contrast, NeuPIMs consumes even higher power (Ours $+$ 760\,W host GPU overhead) while running 2--3.6$\times$ slower due to serialized bank access, yielding 1.8--5.1$\times$ lower Token/J than our design.
These results confirm that near-memory computing alone is insufficient without eliminating the host GPU power tax---our fully integrated 16-NPU architecture achieves both high bandwidth utilization (87--98\% HBM) and high energy efficiency without external GPU dependency.


Fig.~\ref{fig:power_efficiency} shows that our architecture achieves 5.6--6.6$\times$ higher energy efficiency (Token/J) than H100 at BS=1, consistently across Qwen3-235B and Llama4-Maverick.
This advantage holds across batch sizes (BS=1--32) and sequence lengths (4K--1M), driven by the 40\,TB/s aggregate HBM bandwidth of 16~NPUs versus 3.35\,TB/s on H100.
Compared to Rubin, our design maintains a stable 2.6--3.1$\times$ Token/J advantage; against Rubin~TP2, the gap ranges from 2.8$\times$ (1M) to 4.8$\times$ (4K), as TP2's doubled static power and NVLink overhead erode its latency benefit at short sequences.

As shown in Fig.~\ref{fig:power_breakdown}, although our architecture draws higher instantaneous power than H100 (e.g., 1\,316\,W vs.\ 603\,W at BS=4, seq=4K on Llama4), it completes each token 12--16$\times$ faster, resulting in 82--85\% lower energy per token.
In contrast, NeuPIMs consumes even higher power (Ours $+$ 760\,W host GPU overhead) while running 2--3.6$\times$ slower due to serialized bank access, yielding 1.8--5.1$\times$ lower Token/J than our design.
These results confirm that near-memory computing alone is insufficient without eliminating the host GPU power tax---our fully integrated 16-NPU architecture achieves both high bandwidth utilization (87--98\% HBM) and high energy efficiency without external GPU dependency.


Fig.~\ref{fig:power_efficiency} shows that our architecture achieves 5.6--6.6$\times$ higher energy efficiency (Token/J) than H100 at BS=1, consistently across Qwen3-235B and Llama4-Maverick.
This advantage holds across batch sizes (BS=1--32) and sequence lengths (4K--1M), driven by the 40\,TB/s aggregate HBM bandwidth of 16~NPUs versus 3.35\,TB/s on H100.
Compared to Rubin, our design maintains a stable 2.6--3.1$\times$ Token/J advantage; against Rubin~TP2, the gap ranges from 2.8$\times$ (1M) to 4.8$\times$ (4K), as TP2's doubled static power and NVLink overhead erode its latency benefit at short sequences.

As shown in Fig.~\ref{fig:power_breakdown}, although our architecture draws higher instantaneous power than H100 (e.g., 1\,316\,W vs.\ 603\,W at BS=4, seq=4K on Llama4), it completes each token 12--16$\times$ faster, resulting in 82--85\% lower energy per token.
In contrast, NeuPIMs consumes even higher power (Ours $+$ 760\,W host GPU overhead) while running 2--3.6$\times$ slower due to serialized bank access, yielding 1.8--5.1$\times$ lower Token/J than our design.
These results confirm that near-memory computing alone is insufficient without eliminating the host GPU power tax---our fully integrated 16-NPU architecture achieves both high bandwidth utilization (87--98\% HBM) and high energy efficiency without external GPU dependency.


Fig.~\ref{fig:power_efficiency} shows that our architecture achieves 5.6--6.6$\times$ higher energy efficiency (Token/J) than H100 at BS=1, consistently across Qwen3-235B and Llama4-Maverick.
This advantage holds across batch sizes (BS=1--32) and sequence lengths (4K--1M), driven by the 40\,TB/s aggregate HBM bandwidth of 16~NPUs versus 3.35\,TB/s on H100.
Compared to Rubin, our design maintains a stable 2.6--3.1$\times$ Token/J advantage; against Rubin~TP2, the gap ranges from 2.8$\times$ (1M) to 4.8$\times$ (4K), as TP2's doubled static power and NVLink overhead erode its latency benefit at short sequences.

As shown in Fig.~\ref{fig:power_breakdown}, although our architecture draws higher instantaneous power than H100 (e.g., 1\,316\,W vs.\ 603\,W at BS=4, seq=4K on Llama4), it completes each token 12--16$\times$ faster, resulting in 82--85\% lower energy per token.
In contrast, NeuPIMs consumes even higher power (Ours $+$ 760\,W host GPU overhead) while running 2--3.6$\times$ slower due to serialized bank access, yielding 1.8--5.1$\times$ lower Token/J than our design.
These results confirm that near-memory computing alone is insufficient without eliminating the host GPU power tax---our fully integrated 16-NPU architecture achieves both high bandwidth utilization (87--98\% HBM) and high energy efficiency without external GPU dependency.
